\documentclass[../../../include/open-logic-section]{subfiles}

\begin{document}

\olfileid{his}{set}{mythology}

\olsection{أسطورة أكثر منها تاريخًا؟}

وإذ ننظر اليوم في هذه الحوادث، وقد مضى عليها فوق قرن، وجب أن نتثبت قبل
التسليم بما قيل فيها. لا ريب أن النتائج كانت جديدة مثيرة للعجب؛ ولكن هل
بلغت من الصدمة ما يُروى؟ وهل دلّت حقاً على وجوب طرح الحدس الهندسي؟

فأما دعوى الصدمة، فينبّه~\citet{Gouvea2011} إلى أن كلمة كانتور المشهورة
لديدكند، «\emph{je le vois, mais je ne le crois pas}»، انتُزعت في معظم
الروايات من سياقها. وهذا قدر أوسع من ذلك السياق (نقلاً عن
\citeauthor{Gouvea2011}):
\begin{quote}
أرجو أن تعذر حماستي لهذا الموضوع إذا أثقلت عليك بكثرة ما أطلبه من
لطفك وصبرك؛ فالرسائل التي بعثت بها إليك مؤخرًا غير متوقعة وجديدة حتى
بالنسبة إليّ، إلى حد أنني لا أستطيع الاطمئنان حتى أحصل منك، يا صديقي
المكرم، على حكم في صحتها. وما دمت لم توافقني، فلا يسعني إلا أن أقول:
\emph{je le vois, mais je ne le crois pas.}
\end{quote}
فكان كانتور عالماً بأن نتيجته «غير متوقعة وجديدة إلى هذا الحد»، أما أنه
عدّها قط \emph{مما لا يُصدق} ففيه نظر. وإنما كان، كما يبين
\citeauthor{Gouvea2011}، يستثبت ديدكند من صحة البرهان الذي بعث به إليه.

وأما الحدس الهندسي، فإن بيانو نشر منحناه المالئ للفضاء من غير رسم؛ لكن
هيلبرت، حين نشر منحناه، صرح بغرضه: أن يوضح للقراء سبيل فهم نتيجة بيانو إذا
«استعانوا بالحدس الهندسي الآتي». ثم ساق طائفة من \emph{الرسوم} تكاد تطابق
ما ورد في~\olref[his][set][pathology]{sec}.

والأعم من هذا أن الرسوم، وإن قلّ ورودها في البراهين المنشورة، كثيرة الدوران
\emph{جداً} في عمل الرياضيين حين يتوصلون إلى الإثبات. (وتقريب المعنى أن
الحاذقين منهم يعرفون موضع الثقة بالحدس الهندسي.)

وخلاصة الأمر: لا تُسلم للتهويل، أو لا تقبله على الأقل قبل الفحص. ومن رام
البسط فليرجع إلى~\citet{Giaquinto2007}.

\end{document}
